\section{Integrating Inverse Functions}

\begin{enumerate}
  \item Let \(f:I\to J\) be one-to-one. What does it mean for \(f^{-1}:J\to I\) to be the inverse of \(f\)? State both identities involving \(f\) and \(f^{-1}\).
  \item Why is injectivity of \(f\) necessary if we want \(f^{-1}\) to be an ordinary function on \(J=f(I)\)?
  \item Suppose \(x=f(t)\). Explain why
    \[ f^{-1}(x)=t. \]
  \item In the informal substitution
    \[ \int f^{-1}(x)\,dx = \int t f'(t)\,dt, \qquad x=f(t), \]
    which two transformations are being used?
  \item Why should the expression
    \[ dx=f'(t)\,dt \]
    not be regarded as the fundamental rigorous justification of the substitution?
  \item What theorem provides the rigorous justification for the transformation
    \[ \int f^{-1}(x)\,dx \longrightarrow \int f^{-1}(f(t))f'(t)\,dt? \]
  \item Let \(F\) be an antiderivative of \(f^{-1}\). Write explicitly what the statement
    \[ F'=f^{-1} \]
    means.
  \item Differentiate \(F(f(t))\) using the chain rule.
  \item Starting from
    \[ \frac{d}{dt}F(f(t)), \]
    derive in a few lines the identity
    \[ \frac{d}{dt}F(f(t))=t f'(t). \]
  \item Why does the identity
    \[ (F\circ f)'(t)=t f'(t) \]
    show that \(F\circ f\) is an antiderivative of \(t f'(t)\)?
  \item Explain carefully what is meant by the informal identity
    \[ \int f^{-1}(x)\,dx = \int t f'(t)\,dt. \]
    In particular, why are the two antiderivatives being written in different variables?
  \item State a precise definite-integral version of
    \[ \int f^{-1}(x)\,dx = \int t f'(t)\,dt. \]
  \item Prove that for \(a,b\in I\),
    \[ \int_{f(a)}^{f(b)}f^{-1}(x)\,dx = \int_{a}^{b} t f'(t)\,dt. \]
  \item Derive the same definite-integral identity directly from the substitution theorem by setting
    \[ x=f(t). \]
  \item In the substitution
    \[ x=f(t), \]
    what happens to the lower and upper bounds \(f(a)\) and \(f(b)\)?
  \item Does the identity
    \[ \int_{f(a)}^{f(b)}f^{-1}(x)\,dx = \int_{a}^{b} t f'(t)\,dt \]
    require \(f\) to be increasing? Explain.
  \item Suppose \(f\) is decreasing. What is the relative order of \(f(a)\) and \(f(b)\) when \(a<b\), and why does the definite-integral formula still remain valid?
  \item Apply integration by parts to
    \[ \int_{a}^{b} t f'(t)\,dt. \]
    State clearly your choices of \(u\) and \(dv\).
  \item Derive
    \[ \int_{a}^{b} t f'(t)\,dt = b f(b)-a f(a)-\int_{a}^{b} f(t)\,dt. \]
  \item Combine the inverse-function substitution formula with integration by parts to prove
    \[ \int_{f(a)}^{f(b)}f^{-1}(x)\,dx = b f(b)-a f(a)-\int_{a}^{b} f(t)\,dt. \]
  \item Let \(x\in J\) and put
    \[ b=f^{-1}(x). \]
    Show that \(f(b)=x\).
  \item Starting from
    \[ \int_{f(a)}^{f(b)}f^{-1}(u)\,du = b f(b)-a f(a)-\int_{a}^{b} f(t)\,dt, \]
    substitute \(b=f^{-1}(x)\) and derive
    \[ \int_{f(a)}^{x} f^{-1}(u)\,du = x f^{-1}(x)-a f(a) - \int_{a}^{f^{-1}(x)}f(t)\,dt. \]
  \item Why is the previous formula more precise than the informal identity
    \[ \int f^{-1}(x)\,dx = x f^{-1}(x)-\int f(t)\,dt? \]
  \item In the informal formula
    \[ \int f^{-1}(x)\,dx = x f^{-1}(x)-\int f(t)\,dt, \]
    what is ambiguous about the last integral?
  \item Give a corrected indefinite-integral formula in which the dependence on \(x\) is explicit.
  \item Let
    \[ G(x) = x f^{-1}(x) - \int_{a}^{f^{-1}(x)}f(t)\,dt. \]
    Differentiate \(G(x)\) directly and verify that
    \[ G'(x)=f^{-1}(x). \]
  \item In the previous calculation, which rule is needed to differentiate
    \[ \int_{a}^{f^{-1}(x)}f(t)\,dt? \]
  \item Assuming \((f^{-1})'(x)\) exists, carry out the derivative of
    \[ x f^{-1}(x) \]
    using the product rule.
  \item Show explicitly how the terms involving \((f^{-1})'(x)\) cancel when differentiating
    \[ x f^{-1}(x) - \int_{a}^{f^{-1}(x)}f(t)\,dt. \]
  \item Why does the direct differentiation proof of the inverse-function integration formula typically require stronger assumptions than the proof based on
    \[ (F\circ f)'(t)=t f'(t)? \]
  \item What hypotheses are sufficient to ensure that the substitution theorem can be applied to
    \[ \int_{f(a)}^{f(b)}f^{-1}(x)\,dx \]
    with \(x=f(t)\)?
  \item Which assumptions are needed merely to define \(f^{-1}\), and which additional assumptions are needed for the differentiability or integration arguments?
  \item Suppose \(f(t)=t^{2}\) on \([0,\infty)\). Find \(f^{-1}(x)\) and verify
    \[ \int f^{-1}(x)\,dx = \int t f'(t)\,dt \]
    by direct computation.
  \item For \(f(t)=t^{2}\) on \([0,\infty)\), derive the usual formula for
    \[ \int \sqrt{x}\,dx \]
    using the inverse-function method.
  \item Suppose \(f(t)=e^{t}\). Use the formula
    \[ \int f^{-1}(x)\,dx = \int t f'(t)\,dt \]
    to derive an antiderivative of \(\ln x\).
  \item Starting from \(f(t)=e^{t}\), use integration by parts to derive
    \[ \int \ln x\,dx=x\ln x-x+C. \]
  \item Suppose \(f(t)=\tan t\) on \((-\pi/2,\pi/2)\). Use the inverse-function method to express
    \[ \int \arctan x\,dx \]
    as an integral in the variable \(t\).
  \item Complete the derivation of
    \[ \int \arctan x\,dx \]
    from the previous question.
  \item Suppose \(f(t)=\sin t\) on \([-\pi/2,\pi/2]\). Use the inverse-function method to derive a formula for
    \[ \int \arcsin x\,dx. \]
  \item Why must the domain of \(\sin t\) be restricted before treating \(\arcsin\) as its inverse?
  \item Suppose \(f(t)=t^{3}+t\). Without finding an explicit formula for \(f^{-1}\), express
    \[ \int_{f(a)}^{f(b)}f^{-1}(x)\,dx \]
    in terms of \(a\), \(b\), and an elementary integral involving \(f\).
  \item For \(f(t)=t^{3}+t\), compute
    \[ \int_{f(0)}^{f(1)}f^{-1}(x)\,dx \]
    without explicitly solving for \(f^{-1}\).
  \item Explain geometrically why an identity involving
    \[ \int f(x)\,dx \quad\text{and}\quad \int f^{-1}(x)\,dx \]
    should be plausible when \(f\) is increasing.
  \item If the graph of \(y=f(x)\) is reflected across the line \(y=x\), what graph is obtained?
  \item For an increasing positive function \(f\), explain how the rectangle with side lengths \(b\) and \(f(b)\) can be decomposed into areas associated with \(f\) and \(f^{-1}\).
  \item Derive the geometric area identity
    \[ \int_{a}^{b} f(t)\,dt + \int_{f(a)}^{f(b)}f^{-1}(x)\,dx = b f(b)-a f(a). \]
  \item Which part of the rigorous analytic proof corresponds geometrically to the fact that the graphs of \(f\) and \(f^{-1}\) are reflections of each other?
  \item Explain the role of the fundamental theorem of calculus in passing from an antiderivative statement to a definite-integral identity.
  \item Explain the role of the chain rule in the proof.
  \item Explain the role of integration by parts in obtaining a formula involving both \(f\) and \(f^{-1}\).
  \item Distinguish clearly between the following three statements:
    \[ (F\circ f)'(t)=t f'(t), \]
    \[ \int_{f(a)}^{f(b)}f^{-1}(x)\,dx = \int_{a}^{b} t f'(t)\,dt, \]
    and
    \[ \int_{f(a)}^{f(b)}f^{-1}(x)\,dx = b f(b)-a f(a)-\int_{a}^{b} f(t)\,dt. \]
    What theorem or technique is responsible for each transition?
  \item Give a complete proof, from the chain rule and the fundamental theorem of calculus, of
    \[ \int_{f(a)}^{f(b)}f^{-1}(x)\,dx = b f(b)-a f(a)-\int_{a}^{b} f(t)\,dt. \]
  \item Reconstruct the entire argument without using the symbolic differential equation
    \[ dx=f'(t)\,dt. \]
  \item Finally, explain in your own words why the informal substitution
    \[ x=f(t),\qquad dx=f'(t)\,dt \]
    works in this problem, and identify the rigorous theorem that justifies each symbolic step.
\end{enumerate}
